\section{Results and Discussion}
\label{sec:results}
\subsection{Breathing Experience: subjective outcomes}
\label{sec:subjective-results}
The table gives mean scores for control, symmetric visuals and asymmetric visuals, in that order. A positive difference indicates a higher score with immersive visuals. The standardised difference, $d_z$, expresses the change relative to the variability of the paired differences. The paired contrast compares each participant's mean across the two immersive blocks with their control score ($n=39$) [4,5].

\small
\begin{tabular}{p{0.28\linewidth}p{0.28\linewidth}p{0.37\linewidth}}
\toprule
Outcome & Condition means & Immersive conditions versus control \\
\midrule
Breathing synchrony & .787 / .836 / .860 & $t(38)=3.47$, $p=.001$, $d_z=.56$; 95\% confidence interval for the difference [.026, .097] \\
Spatial frame of reference & 3.15 / 4.18 / 4.26 & $t(38)=4.28$, $p<.001$, $d_z=.69$; 95\% confidence interval [.56, 1.57] \\
Perceptual Effects (0--100) & 11.41 / 27.25 / 25.59 & $t(38)=4.90$, $p<.001$, $d_z=.79$ \\
Positive Effects (0--100) & 18.88 / 25.64 / 26.46 & $t(38)=3.22$, $p=.003$, $d_z=.51$ \\
Distressing Effects (0--100) & 12.33 / 17.33 / 17.02 & $t(38)=2.56$, $p=.014$, $d_z=.41$ \\
\bottomrule
\end{tabular}
\normalsize

Immersive feedback increased breathing synchrony and reported spatial self-extension. Participants also reported stronger perceptual and positive experiences and a smaller increase in distress. Of the 11 registered dimensions, six differed after correction for multiple tests: unity, spiritual experience, disembodiment, elementary imagery, audio-visual synaesthesia and changed meaning of percepts. Impaired control/cognition fell short after correction ($p=.038$, $q=.060$), as did anxiety, insightfulness, bliss and complex imagery (unadjusted $p=.063$, .112, .678 and .901).

Body-boundary salience ($p=.864$, $BF_{01}=5.71$) and Small Self ($p=.173$, $BF_{01}=2.39$) showed little evidence of a change. No outcome reliably distinguished symmetric from asymmetric visuals (composite $p\geq.357$). The order-adjusted omnibus tests gave $F(2,75)=19.63$, 7.56 and 4.64 for Perceptual, Positive and Distressing Effects ($q<.001$, .002 and .013). The overall test of breathing synchrony gave $\chi^2(2)=12.66$, $p=.002$ [5].

The predominantly female student sample and exclusion of participants with low synchrony limit generalisation. The design estimates the combined effect of immersive visuals during paced breathing. It cannot separate visual depth, motion and colour, or establish an effect of reversed coupling or time dilation.

\subsection{Respiratory PPS: preliminary behavioural results}
\label{sec:pps-results}
The release contains 42 records and 17,114 trial rows. Cleaning retains 8,916 responses to touch accompanied by sound before linkage exclusions and 8,478 afterwards, leaving 40 behavioural records. Thirty-nine have complete combined behavioural/questionnaire data; questionnaire-only analyses use 41 records [7].

\textbf{Reaction time.} For each person, we subtracted the average touch-only reaction time from the sound-plus-touch reaction time within each visual condition and breathing phase. The 23 October analysis weighted those four combinations equally. All 40 participants contributed data at each sound delay. The Greenhouse--Geisser corrected test is $F(2.52,98.11)=3.28$, $p=.031$, $\eta_p^2=.078$, $\epsilon=.629$ [9]. Corrected means from the shortest to longest sound delay are $-8.60$, $-14.25$, $-22.14$, $-23.87$, $-24.38$ ms. The paired endpoint difference is $-15.78$ ms, 95\% confidence interval [$-31.95$, 0.39]. Including both ambiguous records gives $N=42$, $p=.050$; excluding either alone gives $N=41$, $p=.040$.

Before baseline correction, responses were slower after inhalation than after exhalation. The mixed model estimated a difference of $+17.21$ ms, model-based interval [11.91, 22.50]; a participant-cluster-robust regression gives $+16.02$ ms [5.52, 26.52], unadjusted $p=.00278$. Evidence that the effect of sound delay depended jointly on visual condition and breathing phase was inconclusive ($p=.873$). Temporal expectancy could also account for a change in response speed across sound delays.

\textbf{Boundary identification.} The 17 October analysis fitted an S-shaped curve to estimate the distance at which an approaching sound began to speed responses to touch. Four sets of prior assumptions tested how strongly the data determined the answer. Targets are midpoint in cm and normalized sharpness $1/(4w)$ in cm$^{-1}$. The reference model fits each cell separately. None of 160 non-ambiguous cell midpoint intervals lies wholly within 10--100 cm; only five signed-amplitude intervals exclude zero. Median midpoint interval width is about 173 cm, widening to 259 cm with broader support. Median shape information gain is 0.108 nats.

For distances withheld from fitting, root mean square prediction error was 47.44 ms for the curve and 48.03 ms for a constant response. The difference was $-0.59$ ms, with a paired confidence interval of [$-2.45$, 1.13]. Midpoint contrasts with propagated uncertainty are $-0.53$ [$-23.8$, 23.1] cm for immersive feedback minus control, 0.57 [$-23.2$, 25.3] for phase and $-4.1$ [$-50.5$, 43.5] for interaction. Sharpness intervals span zero. Analyses using only the median boundary estimate discard uncertainty about its location.

\textbf{Sensitivity to cleaning and session time.} In the 20 October joint Bayesian analyses, local trimming retained 8,407 audio-tactile trials. All 48 main parameter intervals included zero. All 30 intervals from the session-time models also included zero, including the 24 from fits meeting convergence criteria. One fit failed that check. Even converged models that required a visible boundary transition predicted the observed profiles poorly [9]. Early-block and time-adjusted analyses could not separate fatigue, practice and carryover.

\subsection{Respiratory PPS: questionnaire findings}
\label{sec:pps-questionnaires}
Of 41 participants with paired ratings, 22 reported greater extension of peripersonal space after Viscereality than after control, six reported less and 13 gave the same rating. The exact paired sign test on the 28 non-tied pairs gave $p=.00372$ [7]. The mean difference was $+0.756$ scale categories, a descriptive summary that assumes equal spacing between categories.

Participants who reported greater trust in their bodily sensations tended to report a larger increase in perceived spatial extent with immersive feedback than with control, Spearman $\rho=.376$, paired-person 95\% interval [.081, .627], unadjusted permutation $p=.0161$. This exploratory comparison was one of eight subscales tested. The decision to report unadjusted values followed inspection of the results. All 92 association intervals and all 192 joint-model coefficient intervals relating questionnaires to behavioural estimates of the boundary included zero under each prior.

\subsection{eMotion: achieved collection and remaining analysis}
\label{sec:s6-results}
Eight participants completed eMotion ($n=8$ of the planned 24), each finishing four conditions under both hand and environment mappings. Three additional participants completed one mapping [10]. Before analysing changes in the time between heartbeats, we must check the quality of each heart recording and match it to the correct condition. Analyses of emotional responses are also pending.

\subsection{Interpretation across experiments}
In both controlled experiments, participants reported that their sense of space extended farther around the body during immersive feedback. In Respiratory PPS, response times varied with sound delay, but the behavioural task did not reliably locate a boundary shift. The clearest evidence concerns the perceived extent of the space around the body.
